\section{Related Work}

\para{Linear Representation Hypothesis (\lrh).} The \lrh suggests that semantic concepts are directions in activation space, allowing for linear manipulation and probing \cite{park2023linear, jiang2024origins}. Early work in word embeddings \cite{elhage2021mathematical} showed that simple vector arithmetic could capture relationships like ``king - man + woman = queen.'' However, recent investigations into the geometry of LLM activations have found that while linear structure is present, it is often non-Euclidean and context-dependent \cite{park2023linear}. Our work extends this by investigating the limits of linear structure when interpolating between \incongruent concepts.

\para{Superposition and Feature Packing.} The phenomenon of superposition \cite{elhage2022toy} explains how models can store more features than they have dimensions by using almost orthogonal directions. \citet{elhage2022toy} used small ReLU networks to demonstrate that features are packed into geometric structures like polytopes. This suggests that while individual features may be linear, their coexistence in a finite-dimensional space creates a complex, sparse landscape. Our work examines how these packed features interact when a forced linear path bridges distant and unrelated regions of this landscape.

\para{Incongruence and Polysemy.} Recent work has explored how models handle multiple semantic meanings simultaneously. \citet{white2022schrodinger} demonstrated that diffusion models can generate visual superpositions of polysemous word senses, showing that the underlying CLIP encodings linearly represent these senses. However, their focus was on visual generation rather than the internal information geometry of the representation space. We build on this by quantifying the ``semantic distance'' between incongruent concepts in probability space.

\para{Sparse Autoencoders (\saes) and Interpretability.} \saes have emerged as a powerful tool for decomposing polysemantic neurons into monosemantic features \cite{cunningham2023sparse, bricken2023monosemanticity}. By training an autoencoder with a sparsity-inducing penalty, researchers can isolate specific semantic components. We use \saes not only for decomposition but as a diagnostic for manifold alignment. If an interpolated activation vector has a high reconstruction error, it indicates that the model has drifted into a region of space that does not correspond to a valid combination of learned features.
